\title{Representation theory of the symmetric groups}
\author{\textit{Based on a lecture series by} Stuart Martin.}
\date{Michaelmas 2026}
\maketitle

\newcommand{\mat}[1]{\left(\begin{smallmatrix} #1 \end{smallmatrix}\right)}
\newcommand{\Mat}[1]{\textnormal{Mat}_{#1 \times #1}}
\newcommand{\adj}{\textnormal{ad}}
\renewcommand{\gg}{\mathfrak{g}}
\newcommand{\bb}{\mathfrak{b}}
\newcommand{\nn}{\mathfrak{n}}
\renewcommand{\gl}{\mathfrak{gl}}
\renewcommand{\sl}{\mathfrak{sl}}
\newcommand{\so}{\mathfrak{so}}
\renewcommand{\sp}{\mathfrak{sp}}
\renewcommand{\sf}{\mathsf}

\newcommand{\GL}[1]{\mathsf{GL}(#1)}

\newcommand{\sym}[1]{\mathfrak{S}_{#1}}

This is a fairly gentle introduction to the representation theory of the
symmetric group. The flavour will be both algebraic and combinatorial, in
particular a first course on representation theory is prerequisite.

Why study representations of symmetric groups? Firstly it is a nice theory. For
instance, there is a \emph{canonical} classification of the irreducibles ---
irreducible representations of $\sym{n}$ over $\bbC$ are parametrised by
\emph{partitions} of $n$. The construction works uniformly for all symmetric
groups, much unlike the case of general finite groups where multiple ad hoc
constructions are needed for each case. 

Another reason is that the theory of symmetric groups serves as fodder for the
theory for other algebras that are diagrammatic in flavour, such as (pure) braid
groups. 

\paragraph{Schur--Weyl duality.} Here is a third reason the representation theory
of the symmetric group is nice. For $V$ a $\bbC$-vector space, the tensor power
$V^{\otimes 2}$ is a representation of $\mathsf{GL}(V)\times \sym{2}$ where
$\mathsf{GL}(V)$ acts on each factor, and $\sym{2}$ acts by permuting factors
(i.e.\ $v\otimes w\mapsto w\otimes v$). Then $V^{\otimes 2}$ splits into
irreducible representations as $\Sym^2 V \oplus \Exter^2 V$; there exists a
correspondence between representations of $\sym{2}$ (namely the trivial and sign
representation) and those of $\mathsf{GL}(V)$ (the two factors of $V^{\otimes
2}$ respectively). \emph{Schur--Weyl duality} extends this theory from $\sym{2}$
to $\sym{r}$. That is, consider the representation $V^{\otimes r}$ of
$\GL{V}\times \sym{r}$ where $\GL{V}$ acts on each factor and $\sym{r}$ permutes
the factors. Then Schur--Weyl duality shows $V^{\otimes r}=\bigoplus_i
A_i\otimes B_i$ where $\{A_i\}$ are irreducible representations of $\GL{V}$ and
$\{B_i\}$ are irreducible representations of $\sym{r}$. Again, there is a
correspondence $A_i\leftrightarrow B_i$ between some representations of $\GL{V}$
and some representations of $\sym{r}$.

Alas this correspondence is not usually bijective --- one issue is that $\GL{V}$
has infinitely many irreducible representations while $\sym{r}$ has a finite
number has a finite number. Another problem is taking inverse powers of
determinants. 

But the correspondence is \emph{almost} bijective if we stabilise ---  if we let
$\dim(V)$ be large then we recover all representations of $\sym{r}$, while if we
let $r$ be large then we recover all representations of $\GL{V}$ up to powers of
the determinant. We won't look at this.

\paragraph{Topology.} A fourth reason comes by taking any topological space $X$
and considering its \emph{configuration space} $\sf{C}^nX$ (i.e.\ the moduli of
$n$ distinct unordered points in $X$). The group $\sym{n}$ naturally acts on
$\sf{C}^n X$, giving a representation of $\sym{n}$ on $\HH^i(\sf{C}^nX, \bbC)$.

\paragraph{In positive characteristic.} Representations over fields of
prime characteristic can be constructed in a `similar' manner to the
characteristic zero theory, which is again unlike the general case for finite
groups (modular representation theory tends to be much, much harder in
prime characteristic than over $\bbC$; even constructing characters is hard).

\subsection{About} This course on symmetric groups is a part of Part III of the
Mathematics tripos at the University of Cambridge. The lectures were delivered in-person, and were transcribed live. The notes can be found
online at \url{https://pas201.user.srcf.net/documents/2026-symmetric-group.pdf}. 

Errors and corrections should be communicated to by email to
\url{pas201@srcf.net}.

\section{Representations and modules}
\marginpar{Lecture 1, 12/10/26} 
We go over some standard results in the representation theory of a finite group
$G$ over a field $\bbk$. By \emph{module}, we will always mean finite
dimensional left modules.

\begin{definition} 
    A \emph{representation} of $G$ on a finite dimensional $\bbk$-vector space
    $V$ is any (and all) of the following.
    \begin{enumerate}[(i)]
        \item a homomorphism $\rho\colon G\to \GL{V}$;
        \item a linear action of $G$ on $V$, i.e.\ a map $G\times V \to V$ such
            that for all $g,h\in G$ and $v,w\in V$ we have $(gh)(v)=g(h(v))$,
            $1(v)=v$, $g(\alpha v + \beta w)=\alpha(gv)+\beta(gw)$;
        \item a $\bbk G$-module structure on $V$, where $\bbk G$ is the
            \emph{group algebra} (the free $\bbk$-module with basis $G$, with
            product given by extending the multiplication in $G$ linearly). 
    \end{enumerate}
\end{definition}

\begin{example}
    The one-dimensional trivial representation of $G$, denoted $\mathbf{1}_G$, is
    given by $V=\bbk$ with $g(v)=v$ for all $g\in G$, $v\in V$.
\end{example}

If $W$ is a subrepresentation of $V$ (i.e.\ a $\bbk G$-submodule) then we write
$W\leq V$. A representation $V$ is \emph{simple} (or \emph{irreducible}) if it
is a simple $\bbk G$-module, i.e.\ it is non-zero and its only
subrepresentations are $\mathbf{0}$ and $V$. We say $V$ is \emph{semisimple} (or
\emph{complete reducible}) if it is a direct sum of simple representations (this
property is closed under taking sub-objects and quotients). We say $V$ is
\emph{indecomposable} if it is not a direct sum of proper subrepresentations, we
say it is \emph{decomposable} otherwise. 

\begin{theorem}[(Maschke)]
    If $\bbk$ has characteristic zero, then every subrepresentation of $V$ is in
    fact a direct summand. In other words, irreducible and indecomposable
    representations of $V$ coincide.
\end{theorem}

A \emph{$G$-homomorphism} between two $\bbk G$-modules $V,W$ is a $\bbk G$-linear
map $\alpha\colon V\to W$.

Suppose $\{L_i\}$ is a collection of pairwise non-isomorphic irreducible
representations of $G$. Suppose $V$ is semisimple. Then there is an isomorphism
$f\colon \bigoplus L_i ^{\oplus n_i}\to V$ for some non-negative integers $n_i$. The
isomorphism is not canonical; but the embedding $f(L_i^{\oplus n_i})\into V$
\emph{is canonical} --- it is the sum of the images of all maps $L_i\to V$,
called the \emph{$L_i$-isotypic piece}. Thus each $\bbk G$-module $V$ is
\emph{canonically} isomorphic to the direct sum of its isotypic components, and
each $L_i$-isotypic component admits a non-canonical decomposition into a direct sum
of $L_i$s.

\begin{theorem}[(Schur's lemma)]
    Suppose $V,W$ are irreducible representations of $G$. Any $G$-homomorphism
    $V\to W$ is either zero or is an isomorphism. In particular, $\End_GV$ is
    a division algebra over $\bbk$ (if $\bbk$ is algebraically closed then
    $\End_GV\simeq \bbk$).
\end{theorem}

\marginpar{Lecture 2, 14/10/26} 
\begin{definition}
    For $G$-representations $L_i, V$ such that $L_i$ is irreducible, the \emph{$L_i$-multiplicity space} in $V$ is $\Hom_G(L_i, V)=M_i$.
\end{definition}

Suppose $\bbk=\ol\bbk$ is algebraically closed, and $V$ is a finite dimensional
isotypic representation of $G$ (that is, $V\simeq L_i^{\oplus n_i}$ for some
irreducible representation $L_i$). In this case, the canonical map $\Hom_G(L_i,
V)\otimes L_i\to V$ is an isomorphism. If $V$ is semisimple and $\{M_i\}$ are the
multiplicity spaces of irreducible representations $\{L_i\}$, then the
analogously defined map $\bigoplus M_i\otimes L_i \to V$ is an isomorphism. If
$V$ is not semisimple, then this map is only injective (and the image of the map
is called the \emph{socle} of $V$).

\begin{proposition}
    The left regular representation of $G$, given by $V=\bbk G$ where $G$ acts by left
    multiplication, is semisimple and has a $(\dim L_i)$-dimensional
    multiplicity space for each irreducible representation $L_i$. In other words
    it decomposes as \[V=\bigoplus_{L_i\in \textnormal{Irr}(G)} L_i^{\oplus \dim L_i}\]
    with corollary that $|G|=\sum (\dim L_i)^2$.
\end{proposition}

\paragraph{Maschke's theorem again.} Call a $\bbk G$-module $P$
\emph{projective} if every short exact sequence of $\bbk G$-modules $0\to N\to
N'\xrightarrow{\alpha} P\to 0$ splits, i.e.\ there exists a map $\beta\colon
P\to N'$ with $\alpha\beta = \id_P$. Every representation of $\bbk G$ is
semisimple if and only if the trivial representation $\bbOne_G$ is projective.
Then Maschke's theorem states that if $G$ is a finite group and the
characteristic of $\bbk$ does not divide $|G|$, then $\bbOne_G$ is projective
and hence every representation of $G$ is semisimple.

\begin{remark}
    If the characteristic of $\bbk$ divides $|G|$, then this \emph{always}
    fails. For example take $G=\bbZ/p$ and $\bbk = \bbF_p$, and consider the
    representation $V=\bbF_p^{\oplus 2}$ where $a\in G$ acts as
    $\mat{1&a\\0&1}$. Then $V$ is a nontrivial self-extension of the trivial
    representation $\bbOne_G$.
\end{remark}

\paragraph{Tensor products.} The \emph{tensor product} $V\otimes_\bbk
W$ of two $G$-representations $V,W$ has the structure of a $G$-module given by
$g(v\otimes w)=(gv) \otimes (gw)$. 

\paragraph{Induction and Restriction.} For $H\leq G$, there is a
\emph{restriction} functor $\text{Res}^G_H\colon \text{Rep}(G)\to \text{Rep}(H)$
where the $\bbk G$ module $V$ is viewed as a $\bbk H$-module in the obvious way.
We sometimes denote the restricted module as $\text{Res}^G_H V =
V\!\downarrow^G_H = V\!\downarrow_H$. It has a left adjoint, the \emph{induction
functor} $\text{Ind}^G_H\colon \text{Rep}(H)\to \text{Rep}(G)$ given by
$\text{Ind}^G_H(-)=\bbk G\otimes_{\bbk H }(-)$. We then have \emph{Frobenius
reciprocity}, \[\Hom_G(\text{Ind}^G_H U, V)= \Hom_H(U, \text{Res}^G_H V).\]

\paragraph{Permutation modules.} A \emph{permutation module} is a representation
of $G$ equipped with a basis on which $G$ acts, i.e.\ a basis $B$ such that
$gb\in B$ for all $g\in G$ and $b\in B$. In this case we denote the module
$\bbk[B]$. For example, the left regular
representation is naturally a permutation module.

\begin{proposition}
    If $G$ acts transitively on a finite set $X$, then the corresponding
    permutation module $\bbk [X]$ is isomorphic to the induced representation
    $\textnormal{Ind}_H^G(\bbOne_H)$ where $H\leq G$ is the stabiliser of any point
    $x\in X$, and $\bbOne_H$ is the trivial representation.
\end{proposition}

\paragraph{Duals.} If $V$ is a finite dimensional
$\bbk$-vector space, then so is $V\dual = \Hom_\bbk (V,\bbk)$. If $e_1,\ldots,
e_n$ is a basis of $V$ then $V\dual$ is naturally equipped with a dual basis
$\varepsilon_1,\ldots, \varepsilon_n$ such that
$\varepsilon_i(e_j)=\delta_{ij}$. If $V$ is a $G$-representation, then so is $V\dual$ via $g\varphi = \varphi \circ g^{-1}$.

\paragraph{Some linear algebra.}
Suppose $V$ is equipped with a symmetric bilinear form $\langle \cdot, \cdot
\rangle$. This induces a natural map $V\to V\dual$ given by $v\mapsto \langle
\cdot, v\rangle$. Even if $\langle \cdot,\cdot \rangle$ is non-singular (i.e.\
$V\orth = 0$)  we could still have a subspace $W\subset V$ with $W\orth \cap W
\neq 0 $. 

So we will be interested in how big such intersections can get. 

Let $e_1,\ldots, e_n$ be a basis of the subspace $W$. The \emph{Gram matrix} of
$W$ with respect to this basis is the $n\times n$ matrix $A$ with $(i,j)$th
entry $\langle e_i, e_j\rangle$.

\begin{proposition}
    The dimension of $W/(W\cap W\orth)$ is the rank of the Gram matrix.

    \begin{proof}
        The induced map $W\to W\dual$ given by $w\mapsto \langle \cdot,
        w\rangle$ is precisely the map given by the Gram matrix when we equip
        $W, W\dual$ with dual bases. The kernel of the map is $W\cap W\orth$, so
        the answer follows from rank--nullity.
    \end{proof}
\end{proposition}

\paragraph{Character theory.} Work over the field $\bbk = \bbC$, and let $V$ be
a finite dimensional $\bbC G$-module. The \emph{character} of $V$ is a function
$\chi_V\colon G\to \bbC$ given by $g\mapsto \text{tr}(g\colon V\to
V)$. This is a \emph{class function}, i.e.\ it is invariant under conjugation.
Some useful calculations --- 
\begin{gather*}
    \chi_{(V\oplus W)} = \chi_V + \chi_W, \qquad \chi_{(V\dual)} = \ol{\chi_V},
    \qquad \chi_{(V\otimes W)} = \chi_V \cdot \chi_W, \qquad \chi_{\Hom(V,W)} =
    \ol{\chi_V}\cdot \chi_W, \\ 
    \text{for the permutation module determined by }G\circlearrowright B, \qquad
    \chi_{\bbk[B]}(g) =\#\{b\in B\;|\; gb = b\}.
\end{gather*}

Character theory is most useful because of the following result. 

\begin{proposition}
    The space of class functions on $G$ has basis given by the characters of
    irreducible representations.
\end{proposition}

In positive characteristic we need to use \emph{Brauer characters}.



